% ------------------------------------------------------------------
% FILE: sections/09_evaluation_framework.tex
% Evaluation framework: metrics, profiling, ablations, and templates.
% \input{} into main.tex — do NOT wrap in \documentclass or \begin{document}.
% ------------------------------------------------------------------

\section{Evaluation Framework}\label{sec:eval}

The algorithm review in \Cref{sec:algorithms} establishes \emph{what}
will be compared; this section defines \emph{how} the comparison will
be conducted.  All metrics, profiling protocols, and ablation
strategies described below constitute a binding evaluation contract:
every candidate model must be assessed against the same criteria to
enable fair ranking.  Crucially, \textbf{all tables in this section
are empty templates}; they will be populated with measured values during
the experimental phase described in \Cref{sec:nextsteps}, not in this
report.


% ------------------------------------------------------------------
\subsection{Primary Metric: Macro-Averaged F1 Score}
\label{subsec:macro_f1}
% ------------------------------------------------------------------

The primary classification metric is the \textbf{macro-averaged F1
score}, defined as the unweighted mean of per-class F1 scores:
%
\begin{equation}\label{eq:macro_f1}
  F_{1}^{\text{macro}}
  = \frac{1}{K} \sum_{k=1}^{K} F_{1,k}
  = \frac{1}{K} \sum_{k=1}^{K}
    \frac{2 \, P_k \, R_k}{P_k + R_k},
\end{equation}
%
where $P_k$ and $R_k$ are the precision and recall for class~$k$,
and $K = 3$~\cite{sokolova2009f1}.

Macro averaging is chosen over accuracy and micro averaging for two
reasons:
%
\begin{enumerate}
  \item \textbf{Class-balance invariance.}  The dataset has mild class
        imbalance ($601 : 501 : 501$, ratio $1.20 : 1 : 1$).  Accuracy
        would overweight performance on the majority class
        (\texttt{dry}); macro~F1 assigns equal importance to all three
        states.
  \item \textbf{Per-class sensitivity.}  A model that sacrifices recall
        on the rarest or most confusable class (potentially
        \texttt{ice}) to boost overall accuracy would be penalised
        under macro~F1 but might appear adequate under accuracy.
\end{enumerate}


% ------------------------------------------------------------------
\subsection{Per-Class Recall and Precision}
\label{subsec:perclass}
% ------------------------------------------------------------------

While macro~F1 is the ranking metric, \textbf{per-class recall}
identifies which physical state is hardest to detect, and
\textbf{per-class precision} identifies which state is most frequently
hallucinated.  Both are reported alongside the confusion matrix for
every model.

\paragraph{Expected confusion patterns.}
Based on the physical similarity of the three states, the most
confusable class pair is expected to be \texttt{wet} vs.\
\texttt{ice}: both involve a thin layer of water-phase material on
the sensor surface, and their $S_{11}$ signatures may differ primarily
in subtle resonance-frequency shifts and magnitude attenuation changes
caused by the dielectric constant difference between liquid water and
ice.  The \texttt{dry} state, representing a bare sensor, is expected
to be the easiest to distinguish.  The confusion matrix will confirm
or refute this hypothesis.


% ------------------------------------------------------------------
\subsection{Confusion Matrix Analysis}
\label{subsec:confusion}
% ------------------------------------------------------------------

For each trained model, the $3 \times 3$ confusion matrix
$\mathbf{C} \in \mathbb{N}^{3 \times 3}$ will be computed on the
held-out test partition, where $C_{ij}$ is the number of samples
with true label~$i$ predicted as label~$j$.  In addition to the raw
counts, the normalised confusion matrix (each row divided by the row
sum) will be reported to enable cross-model visual comparison.

Specific patterns to examine:
%
\begin{itemize}
  \item Off-diagonal concentration: is misclassification symmetric
        (i.e., \texttt{wet}$\to$\texttt{ice} as frequent as
        \texttt{ice}$\to$\texttt{wet})?
  \item Diagonal dominance: does the model achieve near-perfect
        recall on \texttt{dry} while struggling with the other two?
  \item Error magnitude: are errors single-sample noise or
        systematic bias toward one class?
\end{itemize}


% ------------------------------------------------------------------
\subsection{Robustness Testing}
\label{subsec:robustness}
% ------------------------------------------------------------------

Real-world deployment exposes the model to measurement noise, sensor
degradation, and environmental variation not present in laboratory
data.  To estimate how gracefully each model degrades, the evaluation
includes a \textbf{robustness stress test}:
%
\begin{enumerate}
  \item Take the fixed held-out test set (no retraining).
  \item Inject additive white Gaussian noise at calibrated
        signal-to-noise ratios (SNR):
        $\text{SNR} \in \{40, 30, 20, 15, 10\}$~dB.
  \item Re-evaluate macro~F1 at each noise level.
  \item Plot the resulting \textbf{accuracy degradation curve}:
        macro~F1 vs.\ injected SNR.
\end{enumerate}

Models that maintain high macro~F1 at lower SNR are more robust.
A sharp cliff in the degradation curve indicates that the model's
decision boundary is fragile and positioned close to the data
manifold.  A gradual slope indicates robust margin.

For the 2D route (Route~B), an analogous test will apply pixel-level
Gaussian noise to the rasterised images and measure classification
stability.


% ------------------------------------------------------------------
\subsection{Model Size Profiling}
\label{subsec:model_size}
% ------------------------------------------------------------------

Edge deployment imposes strict constraints on model memory footprint.
For each candidate model, the following size metrics will be recorded:
%
\begin{itemize}
  \item \textbf{Total trainable parameters:} the number of scalar
        values that define the model (weights and biases).
  \item \textbf{Model file size:} the serialised model size on disk
        in float32, float16, and INT8 formats.
  \item \textbf{Peak inference memory:} the maximum RAM required during
        a single forward pass, including intermediate activations.
  \item \textbf{Estimated FLOPS:} the number of floating-point
        multiply-accumulate operations for one forward pass, computed
        analytically or via profiling tools.
\end{itemize}

These values determine whether a model can fit within the memory budget
of the target edge hardware class (see \Cref{sec:nextsteps}).  Models
that exceed the budget after INT8 quantisation will be flagged for
pruning or architectural reduction before re-evaluation.


% ------------------------------------------------------------------
\subsection{Inference Latency Profiling}
\label{subsec:latency}
% ------------------------------------------------------------------

Classification must complete within a real-time budget determined by
the application's sensing cadence.  Inference latency will be measured
as:
%
\begin{itemize}
  \item \textbf{Wall-clock time per sample} on desktop hardware
        (CPU-only, no GPU) as a normalised comparison baseline.
  \item \textbf{Wall-clock time per sample} on representative edge
        hardware (ARM Cortex-A class processor), where available.
  \item \textbf{Throughput} (samples per second) to assess batch
        processing capability.
\end{itemize}

All latency measurements will be averaged over at least \num{100}
inference calls after a warm-up phase to exclude JIT compilation and
cache-warming effects.  Standard deviation will be reported alongside
the mean.


% ------------------------------------------------------------------
\subsection{Calibration Analysis (Optional, Later Stage)}
\label{subsec:calibration}
% ------------------------------------------------------------------

Well-calibrated models produce confidence scores that accurately
reflect the true probability of correctness~\cite{guo2017calibration}.
This property is particularly valuable for edge sensing, where a
high-confidence misclassification (e.g., confidently labelling
\texttt{ice} as \texttt{dry}) can have operational consequences.

If calibration analysis is conducted, it will include:
%
\begin{itemize}
  \item \textbf{Reliability diagrams:} plot predicted confidence
        (binned) against observed accuracy per bin.
  \item \textbf{Expected Calibration Error (ECE):} the weighted
        average of the absolute difference between confidence and
        accuracy across bins:
        %
        \begin{equation}\label{eq:ece}
          \text{ECE} = \sum_{m=1}^{M} \frac{|B_m|}{n}
          \bigl| \mathrm{acc}(B_m) - \mathrm{conf}(B_m) \bigr|,
        \end{equation}
        %
        where $B_m$ is the $m$-th confidence bin, $|B_m|$ the number of
        samples in it, and $n$ the total sample count.
  \item \textbf{Post-hoc recalibration:} temperature scaling or Platt
        scaling applied to the model's logits, trained on a held-out
        calibration set.
\end{itemize}

Calibration analysis is marked as optional because it is a post-training
diagnostic; it does not affect model selection but may influence the
deployment configuration.


% ------------------------------------------------------------------
\subsection{Abstention and Reject Option}
\label{subsec:abstention}
% ------------------------------------------------------------------

For safety-critical or operationally sensitive applications, a model
should have the option to \textbf{abstain} from prediction when its
confidence is below a threshold, rather than risk a misclassification.
The abstention protocol operates as follows:
%
\begin{enumerate}
  \item Define a confidence threshold $\tau \in (0, 1)$.
  \item For each test sample, compute the maximum predicted class
        probability $p_{\max} = \max_k P(y = k \mid \mathbf{x})$.
  \item If $p_{\max} < \tau$, the model abstains (outputs ``uncertain'').
  \item Report:
        \begin{itemize}
          \item \textbf{Coverage}: fraction of test samples where the
                model makes a prediction ($p_{\max} \geq \tau$).
          \item \textbf{Selective accuracy}: accuracy only on
                non-abstained samples.
          \item \textbf{Coverage--accuracy trade-off curve}:
                sweep $\tau$ from 0 to 1 and plot coverage vs.\
                selective accuracy.
        \end{itemize}
\end{enumerate}

This mechanism is complementary to calibration: a well-calibrated model
will produce a smooth, monotonically increasing selective-accuracy
curve as $\tau$ increases.  A poorly calibrated model will show erratic
behaviour.


% ------------------------------------------------------------------
\subsection{Representation-Specific Ablations}
\label{subsec:ablations}
% ------------------------------------------------------------------

The representation routes defined in \Cref{sec:repr} offer multiple
encoding choices.  The evaluation framework requires systematic
ablation studies to determine which encoding yields the best
classification performance.

\paragraph{Route~A ablations.}
Each supervised model will be trained and evaluated on the following
1D encodings:
%
\begin{enumerate}
  \item Magnitude only ($|S_{11}|$ in dB, $d = 4001$).
  \item Phase only ($\angle S_{11}$ in degrees, $d = 4001$).
  \item Magnitude + phase (dual-channel, $d = 8002$ or
        $L = 4001 \times C = 2$).
  \item Magnitude + unwrapped phase (dual-channel).
  \item Magnitude + first derivative of magnitude
        ($d = 4001 + 4000 = 7{,}001$).
\end{enumerate}

\paragraph{Route~B ablations.}
For models compatible with 2D inputs, the following rasterisation
parameters will be varied:
%
\begin{enumerate}
  \item Image resolution: $64 \times 64$, $128 \times 128$,
        $224 \times 224$.
  \item Channel configuration: single-channel (magnitude only),
        dual-channel (magnitude + phase).
  \item Rasterisation method: direct curve plot vs.\ heatmap
        interpolation.
\end{enumerate}

\paragraph{Ablation reporting.}
Each ablation varies exactly one encoding dimension while holding all
other pipeline elements (preprocessing, normalisation, train--test
split, hyperparameters) constant.  Results will be reported in a
factorial table with models as rows and encodings as columns.


% ------------------------------------------------------------------
\subsection{Supervised vs.\ Unsupervised Evaluation Boundary}
\label{subsec:sup_vs_unsup}
% ------------------------------------------------------------------

K-Means (\Cref{subsec:kmeans}) is the only unsupervised method in
the comparison.  Its evaluation follows a fundamentally different
protocol from the nine supervised models:
%
\begin{itemize}
  \item K-Means output is a set of cluster assignments, not class
        predictions.  Cluster-to-class mapping is performed via
        Hungarian matching against the true labels.
  \item The primary metric for K-Means is the \textbf{Adjusted Rand
        Index (ARI)}, not macro~F1.
  \item Silhouette scores measure internal cluster quality
        independently of labels.
  \item K-Means results \emph{inform} but \emph{do not replace}
        supervised evaluation.  A high ARI provides evidence that the
        feature space is well-structured; it does not establish that
        a nearest-centroid classifier is the best deployment choice.
\end{itemize}

K-Means results will be reported separately from the supervised model
ranking and will be interpreted as a diagnostic of representation
quality rather than a competing classification method.


% ------------------------------------------------------------------
\subsection{Evaluation Table Templates}
\label{subsec:eval_tables}
% ------------------------------------------------------------------

The following tables define the structure that experimental results
will follow.  \textbf{All cells are intentionally left empty}; they
will be populated during the experimental phase.

% --- Template 1: Macro F1 comparison ---
\begin{table}[htbp]
  \centering
  \caption{Macro-averaged F1 scores across models and Route~A encodings.
           Values to be filled during the experimental phase.}
  \label{tab:macro_f1_results}
  \begin{tabular}{l c c c c c}
    \toprule
    \textbf{Model}
      & \textbf{Mag}
      & \textbf{Phase}
      & \textbf{Mag+Phase}
      & \textbf{Mag+Unwrap}
      & \textbf{Mag+Deriv} \\
    \midrule
    Logistic Regression   & --- & --- & --- & --- & --- \\
    Linear SVM            & --- & --- & --- & --- & --- \\
    RBF SVM               & --- & --- & --- & --- & --- \\
    Random Forest         & --- & --- & --- & --- & --- \\
    Gradient Boosted Trees & --- & --- & --- & --- & --- \\
    Multilayer Perceptron & --- & --- & --- & --- & --- \\
    1D CNN                & --- & --- & --- & --- & --- \\
    ResNet-1D             & --- & --- & --- & --- & --- \\
    KAN                   & --- & --- & --- & --- & --- \\
    \bottomrule
  \end{tabular}
\end{table}

% --- Template 2: Per-class recall ---
\begin{table}[htbp]
  \centering
  \caption{Per-class recall for each model (best-performing encoding).
           Values to be filled during the experimental phase.}
  \label{tab:perclass_recall}
  \begin{tabular}{l c c c}
    \toprule
    \textbf{Model}
      & \textbf{Recall (\texttt{dry})}
      & \textbf{Recall (\texttt{wet})}
      & \textbf{Recall (\texttt{ice})} \\
    \midrule
    Logistic Regression    & --- & --- & --- \\
    Linear SVM             & --- & --- & --- \\
    RBF SVM                & --- & --- & --- \\
    Random Forest          & --- & --- & --- \\
    Gradient Boosted Trees & --- & --- & --- \\
    Multilayer Perceptron  & --- & --- & --- \\
    1D CNN                 & --- & --- & --- \\
    ResNet-1D              & --- & --- & --- \\
    KAN                    & --- & --- & --- \\
    \bottomrule
  \end{tabular}
\end{table}

% --- Template 3: Model complexity ---
\begin{table}[htbp]
  \centering
  \caption{Model complexity and edge deployment metrics.
           Values to be filled during the experimental phase.}
  \label{tab:model_complexity}
  \begin{tabular}{l r r r r}
    \toprule
    \textbf{Model}
      & \textbf{Parameters}
      & \textbf{Size (INT8)}
      & \textbf{FLOPS}
      & \textbf{Latency} \\
    \midrule
    Logistic Regression    & --- & --- & --- & --- \\
    Linear SVM             & --- & --- & --- & --- \\
    RBF SVM                & --- & --- & --- & --- \\
    Random Forest          & --- & --- & --- & --- \\
    Gradient Boosted Trees & --- & --- & --- & --- \\
    Multilayer Perceptron  & --- & --- & --- & --- \\
    1D CNN                 & --- & --- & --- & --- \\
    ResNet-1D              & --- & --- & --- & --- \\
    KAN                    & --- & --- & --- & --- \\
    \bottomrule
  \end{tabular}
\end{table}

\paragraph{Experimental-phase contract.}
These templates encode the evaluation contract.  During the
experimental phase, every supervised model must appear in every row,
and every encoding variant must appear in every column.  Models that
cannot be trained on a particular encoding (e.g., if Route~B is
incompatible) will be marked as ``N/A'' rather than omitted.  No
partial or selective reporting is acceptable.
